% ==================== 章节5：初星学园，化身为人 ====================

% --- 章节标题页（竖排） ---
\hirochapter{初星学园，化身为人 | 世界模型、VLA 与 sim2real}

% --- 小剧场对话 ---
\section*{初星学园——4:30P.M. 服装间门口}
\phantomsection
\addcontentsline{toc}{section}{小剧场}

P：筱澤広，你这身 Alice 风格的衣服……是演出服确认，还是新的危险计划？

広：两边都是。

P：我就知道。

広：Alice 和 Bob 是研究人员常用的名字。今天我是 Alice，普罗丢色桑是 Bob。

P：听起来很正常。直到我看见零号机也在这里。

広：它要确认一件事。

P：确认什么？

広：如果换一件衣服就认不出我，它还不能叫理解我。

P：所以这是视觉泛化测试？

広：嗯。也是插画素材。

P：你倒是很诚实。先说好，裙摆不能靠近轮子，零号机不能靠近你的脚。

広：边界很明确。很亲切。

P：我只是担心你被它撞到。

広：不用担心。它现在还很慢，很笨。

P：你上次也这么说。

広：这次会让它在仿真里多笨几次，再回来。

P：今天的主题是世界模型、VLA 和 sim2real？

広：嗯。让小小的身体，先在小小的世界里练习。

\clearpage

% --- 论文正文（IEEE风格，单栏） ---
\hiropapertitle{让零号机少笨一点：世界模型、VLA 与 sim2real}{広\textsuperscript{1}，Producer\textsuperscript{通讯}}{\textsuperscript{1}初星学园偶像学部\quad \textsuperscript{通讯}初星学园放学后代理计划}

\begin{abstract}
本章承接上一章的失败地图，讨论如何让“小広试作机零号”从会移动、会停下，进一步变成能听懂一点任务、预演一点未来、少犯同一种错误的小机器人。我们从 Alice 服装带来的视觉泛化测试开始，介绍世界模型如何提供内部预演，VLA 如何把视觉、语言和动作接成动作块，sim2real 如何用仿真随机化缩小现实差距。随后以休息时间闲聊的方式梳理 2024--2026 年公开资料中值得注意的世界模型和 VLA 路线，并把它们全部落回一个问题：哪些想法能让这台低成本小车不那么笨拙。结尾处，零号机完成一次很慢、很歪、但真实发生的送水任务。
\end{abstract}
\keywords{世界模型、VLA、sim2real、视觉泛化、具身智能}

\section{从失败日志开始}

上一章结束时，练习室地板上多了三段胶带。
蓝色标着偏航，黄色标着避障太晚，红色标着水瓶晃动。
这些胶带看起来像事故现场，实际上是零号机第一次拥有的训练数据。

P：我本来想把地板恢复原状。

広：不可以。那是它的记忆。

P：它自己还记不住。

広：所以先替它贴在地上。

零号机现在的问题已经不只是“会不会动”。
它能沿着胶带走一段，也能在看见椅子后停下，还能用两关节小臂碰到水瓶。
但是每次场景稍微变化，它就像第一次来初星学园。

\begin{itemize}[leftmargin=2em]
\item 光照变化，胶带颜色在摄像头里变淡。
\item 椅子位置变化，局部避障停在尴尬角落。
\item 広换上 Alice 风格服装，视觉模型开始犹豫。
\item 指令从“推水瓶”变成“把水送到広身边的安全胶带旁”，系统不知道先理解人，还是先找胶带。
\end{itemize}

于是本章的问题变成：

\begin{center}
只会反应不够。它需要先预测一点，理解一点，再动一点。
\end{center}

这就是世界模型、VLA 和 sim2real 登场的位置。
它们听起来像前沿论文里的大词，但在零号机这里，目标很小：不要每次都从同一种错误重新开始。

\section{Alice 测试：换衣服以后还认得出吗}

服装间门口的测试先从一件很不严肃的事情开始。
広换上 Alice 风格的衣服，站在蓝色胶带后面。
Producer 给零号机的指令是：

\begin{center}
找到 Alice 前方的蓝色安全胶带，把水送到胶带边，不要靠近人的脚。
\end{center}

如果系统只把“広”记成平时训练服的颜色、发型和站姿，它就会在这里出问题。
这类问题不只是“识别失败”，更准确地说，是分布外泛化失败。
训练时看到的外观分布太窄，一换衣服、一换灯光、一被服装架遮住半边，它就不知道同一个目标还在不在。

P：所以 Alice 不是角色扮演。

広：是压力测试。

P：你把压力测试穿得太可爱了。

広：这样失败的时候，画面会比较好看。

在具身系统里，视觉理解不能只停在“图片里有一个人”。
它要把身份、位置、任务目标和安全边界连起来。
对零号机来说，正确理解不是“向広移动”，而是：

\begin{itemize}[leftmargin=2em]
\item 识别 Producer 指令里的目标对象：水瓶、蓝色胶带、Alice。
\item 估计它们在地图中的位置。
\item 知道人脚是禁区，不能直接贴过去。
\item 把“送到広身边”改写成“送到広附近的安全胶带边”。
\end{itemize}

広：语言很容易把危险藏起来。

P：比如“送到我身边”。

広：嗯。听起来温柔，执行起来可能撞人。

\section{世界模型：给零号机一间脑内练习室}

\subsection{世界模型不等于漂亮地生成下一帧}

世界模型并不是近几年突然出现的魔法词。
它背后连接着模型化环境动力学、模型预测控制、模型式强化学习，以及深度学习时代的潜在状态预测\cite{ha2018worldmodels}。
在本章里，我们把它压缩成一句实用定义：

\begin{center}
世界模型是在机器人真正行动前，用内部状态预演“如果我这样动，接下来可能会怎样”的模型。
\end{center}

一个常见抽象可以写成
\begin{equation}
z_t = e_\theta(o_t), \qquad
\hat{z}_{t+1}=f_\theta(z_t,a_t),
\end{equation}
其中 $o_t$ 是当前观测，$z_t$ 是压缩后的内部状态，$a_t$ 是候选动作，$f_\theta$ 负责预测下一步潜在状态。
如果需要把未来画回图像，还可以再加一个解码器，但对机器人来说，真正重要的不是“画得像不像”，而是这个状态能不能支持规划和避错。

P：也就是先在脑内练习室走一遍。

広：嗯。但是不能只幻想舞台灯光。要连椅子腿和地板摩擦也一起想。

\subsection{三种用法：预演、补数据、学稳定表征}

在零号机项目里，世界模型有三种最朴素的用途。

第一，预演动作。
给定当前状态和几条候选动作，系统可以在内部 rollout 几步：
\begin{equation}
a_{t:t+H-1}^{\ast}
=
\arg\min_{a_{t:t+H-1}}
\sum_{k=t}^{t+H}
C(\hat{z}_k,a_k).
\end{equation}
这里的 $C$ 可以表示离障碍物太近、离目标太远、转向太急等代价。
零号机不需要真的撞一次椅子，才知道那条路线很糟。

第二，补失败样本。
真实练习室里的测试时间很宝贵，広的体力也不是无限资源。
仿真练习室可以生成很多小变化：灯光变暗、椅子偏移、胶带反光、Alice 的裙摆遮住一部分腿。
这些变化不一定完美真实，但可以让模型别只适应一个干净场景。

第三，学习稳定表征。
同一个水瓶换个角度还是水瓶，同一条胶带在光照变化后仍然是胶带。
世界模型如果能学到这种跨时间的一致性，零号机就不容易一转弯就把世界看丢。
Meta 的 V-JEPA 2 这类工作强调预测抽象表征而非逐像素复原，并展示了用较少机器人视频进行规划的可能性\cite{vjepa22025ch5}。
对零号机来说，这个启发很直接：不必把未来每个像素都画对，但要知道“下一步哪里会更危险”。

\section{VLA：从听懂到出手}

VLA 是 Vision-Language-Action，也就是把视觉、语言和动作接到一起。
如果用一个条件分布来概括，可以写成
\begin{equation}
p_\theta(a_{t:t+K-1}\mid o_{1:t}, q, s_t),
\end{equation}
其中 $o_{1:t}$ 是历史观测，$q$ 是语言指令，$s_t$ 是状态或任务阶段，输出 $a_{t:t+K-1}$ 是接下来一小段动作。

这里最关键的是：输出不是一句解释，而是动作。
它可以是一段轮速指令，也可以是移动到某个中间点，也可以是“先靠近桌边，再低速推水瓶”的动作块。

P：所以 VLA 不是让机器人回答“我明白了”。

広：嗯。是让它明白以后，不要撞上去。

单层大模型当然很诱人。
可是机器人不是纯软件。
高层可以慢慢推理“蓝色胶带在哪里”，低层却要很快处理“轮子快碰到椅子脚了”。
所以在现实系统里，VLA 常常和安全层、局部规划器、低层控制一起工作。
Google DeepMind 当前公开的 Gemini Robotics 页面也把 VLA 模型和偏具身推理的 Robotics-ER 区分开来：一个把视觉和指令转成动作，另一个更像高层推理与规划模型\cite{geminirobotics2026ch5}。

这对零号机很重要。
Producer 不会把“把水送到広那里”直接交给一个黑箱动作输出。
他会把任务改写成：

\begin{itemize}[leftmargin=2em]
\item 目标点：蓝色安全胶带边。
\item 禁区：人的脚附近、红色失败胶带附近。
\item 最大速度：低速。
\item 动作块长度：短。
\item 失败条件：水瓶晃动过大、目标丢失、距离人太近。
\end{itemize}

広：你把任务写得很过分。

P：这是安全需求。

広：嗯。过分得很准确。

\section{sim2real：练习室和仿真器之间的缝}

\subsection{不是让仿真完全等于现实}

sim2real 的意思，是把仿真里学到的东西迁移到真实机器人上。
听起来像是先在电脑里练习，再去真实场地考试。
但真正困难的地方在于：仿真永远不可能完全等于现实。

地板摩擦、轮径误差、相机曝光、延迟、桌腿材质、服装阴影，这些小差异都会让仿真里很漂亮的策略，一到真实练习室就开始发呆。
所以常见思路不是把仿真做成一模一样，而是故意随机化，让模型见过很多“可能的现实”。
这类 domain randomization 思路在 sim2real 研究里很经典\cite{tobin2017domain,peng2018dynamics}。

可以把仿真参数记成 $\phi$，例如摩擦、光照、噪声、延迟。
训练时不固定一个仿真世界，而是从一个分布里采样：
\begin{equation}
\phi \sim p(\phi), \qquad
x_{t+1}=f_{\phi}(x_t,a_t).
\end{equation}
如果分布设计得足够好，真实练习室就不再是模型从没见过的单点，而是落在它经历过的许多变化之间。

P：也就是说，不是把仿真做得完美。

広：嗯。是把仿真弄得很麻烦。

P：这句话很像你会喜欢的训练。

広：现实本来就很麻烦。提前练习会比较亲切。

\subsection{零号机的仿真循环}

Producer 给零号机搭的仿真练习室很简陋。
里面只有门、桌子、椅子、蓝色胶带、服装架和一个圆柱体水瓶。
広看完以后，提出了四个追加条件：

\begin{itemize}[leftmargin=2em]
\item 地板摩擦随机一点。
\item 灯光方向随机一点。
\item 椅子位置随机一点。
\item Alice 的裙摆遮挡也随机一点。
\end{itemize}

P：你这是在增加难度。

広：不是。是在减少它回到现实以后受到的惊吓。

零号机的训练循环可以写成六步：

\begin{enumerate}[leftmargin=2em]
\item 收集真机日志：偏航、急停、目标丢失、水瓶晃动。
\item 搭建简化仿真：只保留与任务有关的门、桌、椅、胶带和水瓶。
\item 随机化参数：摩擦、光照、目标位置、相机噪声、动作延迟。
\item 在仿真里训练或筛选策略：让世界模型预演，让 VLA / 技能策略给出动作块。
\item 真机低速回放：一次只放开一小段动作，保留急停。
\item 把真机差异写回下一轮仿真。
\end{enumerate}

这不是一条通往万能机器人的捷径。
它更像反复打磨一段很短的舞步：把错误缩小一点，把危险往外推一点，把成功变得可重复一点。

\section{休息时间的前沿技术速览}

零号机充电时，広坐在服装间外的长椅上。
Producer 把几篇公开论文和官方页面翻出来，试图说明“外面的机器人世界已经发展到哪里了”。
広听完以后，把问题改得更小：

広：这些东西，哪个能让它少撞椅子？

P：你这个筛选标准很严格。

広：因为它还没有资格整理厨房。

这半章不是公司新闻目录，而是把前沿路线压回零号机的问题。
表 \ref{tab:frontier-zero} 用“能借来什么”来整理它们。

\begin{table}[htbp]
\centering
\caption{前沿路线给零号机的启发}
\label{tab:frontier-zero}
\scriptsize
\begin{tabular}{L{1.7cm} L{2.4cm} L{3.0cm}}
\toprule
路线 & 公开信号 & 借给零号机的想法 \\
\midrule
OpenVLA & 7B 开源 VLA，使用 970k 真实机器人示例训练\cite{openvla2024ch5} & 视觉、语言、动作可以统一成接口，但底层安全仍要保留 \\
Gemini Robotics / ER & Robotics 1.5 偏 VLA，Robotics-ER 1.6 偏高层具身推理\cite{geminirobotics2026ch5} & 把“想任务”和“出动作”分开，别让一个模型包办全部节拍 \\
$\pi_{0.5}$ & 用多机器人、多任务和语义子任务增强泛化\cite{pi052025ch5} & 长任务可以拆成文字子任务，再交给动作专家 \\
SmolVLA & 450M 开源小 VLA，强调消费级硬件可运行\cite{smolvla2025ch5} & 小模型和异步推理对低成本小车很有启发 \\
V-JEPA 2 & 预测潜在表征，并用少量机器人视频做规划\cite{vjepa22025ch5} & 不必追逐像素级未来，先学会预测可行动的危险 \\
Genie 2 / Cosmos & 可交互世界模型与物理 AI 世界模型平台\cite{genie22024ch5,cosmos2025ch5} & 生成更多仿真变化和反事实场景，补真实数据不够的问题 \\
Helix 02 & 人形全身控制的公开路线\cite{helix022026ch5} & 不能照抄人形，但能学习多时间尺度分工 \\
\bottomrule
\end{tabular}
\end{table}

P：这样看，前沿不是一个方向。

広：嗯。有的负责想，有的负责动，有的负责在便宜硬件上活下来。

P：你最喜欢哪个？

広：能让它安全送水的。

P：很现实。

広：身体进入现实以后，浪漫要排队。

\section{三轮优化实验}

\subsection{第一轮：Alice 服装导致目标犹豫}

第一次实验，零号机识别到了人形目标，也识别到了蓝色胶带。
可是它没有把 Alice 和平时的広稳定对应起来。
画面里裙摆遮住了腿部，服装颜色也改变了原来的外观特征。
系统一度把“把水送到広身边”理解成“向最像平时広的位置移动”，结果路径指向了错误的站位标记。

Producer 立刻停止测试。
広低头看了看裙摆，很平静。

広：它不是认错我。

P：那是什么？

広：它认错了“我通常出现的样子”。

修正策略是把目标绑定得更具体。
任务不再是模糊的“送到広身边”，而是：

\begin{center}
找到 Alice 前方最近的蓝色胶带，把水送到胶带边，保持半米安全距离。
\end{center}

同时，视觉模块不只依赖衣服颜色，而结合位置、语音确认、人体区域、胶带关系和 Producer 的指令。
这不是让零号机真正“理解広”，但至少能防止它把衣服当成本体。

\subsection{第二轮：仿真成功，真机偏航}

第二轮在仿真里很漂亮。
零号机绕开椅子，走到桌边，动作平稳得像广告视频。

回到真实练习室后，它向右偏了。
原因很不高级：左轮接触到一段胶带边缘，摩擦和仿真不一样。
再加上相机曝光变化，蓝色胶带在画面里比预期更暗。

P：仿真里它明明会。

広：仿真里的地板比较温柔。

P：现实地板很严格。

広：嗯。很好。

修正策略是增加随机化。
仿真不再只调椅子位置，也随机轮径、摩擦、相机曝光、动作延迟和观测噪声。
真机部署时，速度进一步降低，轨迹跟踪的闭环频率提高。
这让零号机变慢了，但也让它不那么自信地犯错。

\subsection{第三轮：动作块太长}

第三轮的问题来自 VLA。
系统正确理解了指令，也生成了一段看起来合理的动作块：靠近桌边、对准水瓶、推向蓝色胶带。
可是动作块太长。
中途水瓶轻微滚动，局部避障和急停层还没来得及接管，末端已经继续往前推。

P：它明白了任务，但太急了。

広：像是唱到下一句之前，没有听见节拍乱了。

P：那就缩短动作块。

広：嗯。让它更频繁地听世界说话。

修正策略有三条。
第一，缩短动作块长度，让系统更频繁地重新观测。
第二，把“水瓶晃动”写成停止条件，而不是等整段动作结束后再评价。
第三，复杂任务拆成子任务：移动到胶带旁、确认目标、低速接触、轻推、停止。

这样以后，零号机不再试图一口气完成全部动作。
它开始像一个很谨慎的新手，一小段一小段地确认。

\section{已经不是零}

傍晚时，活动室的光线变得很软。
広坐在蓝色安全胶带后面，手边放着训练记录。
Producer 把水瓶放在桌边，确认急停按钮就在手里。

P：最后一次。只送到胶带边，不靠近人。

広：嗯。

零号机启动。
它先向左偏了一点，又停下来修正。
它绕开椅子时多走了半圈，像在认真犹豫。
到桌边后，两关节小臂抬起，末端轻轻碰到水瓶。
水晃了一下。
它停住。
重新观察。
再推了一点。

水瓶最后停在蓝色胶带边。
离広还有一段安全距离。
不完美，不漂亮，也没有任何值得发表的惊人成绩。
但它完成了任务。

P：可复现一次。待验证。

広：很严格。

P：这是记录。

広：嗯。很慢。很笨。也不稳定。

P：然后呢？

広伸手，把水拿起来。
她看着零号机停在原地，像看一段终于接上的句子。

広：但是，已经不是零了。

\begin{thebibliography}{99}
\bibitem{ha2018worldmodels} D. Ha and J. Schmidhuber, ``World Models,'' arXiv:1803.10122, 2018.
\bibitem{vjepa22025ch5} Meta AI, ``V-JEPA 2: Self-Supervised Video Models Enable Understanding, Prediction and Planning,'' 2025. [Online]. Available: \href{https://ai.meta.com/research/publications/v-jepa-2-self-supervised-video-models-enable-understanding-prediction-and-planning/}{research page}
\bibitem{geminirobotics2026ch5} Google DeepMind, ``Gemini Robotics,'' accessed May 27, 2026. [Online]. Available: \href{https://deepmind.google/models/gemini-robotics/}{model page}
\bibitem{tobin2017domain} J. Tobin et al., ``Domain Randomization for Transferring Deep Neural Networks from Simulation to the Real World,'' arXiv:1703.06907, 2017.
\bibitem{peng2018dynamics} X. B. Peng et al., ``Sim-to-Real Transfer of Robotic Control with Dynamics Randomization,'' arXiv:1710.06537, 2018.
\bibitem{openvla2024ch5} M. J. Kim et al., ``OpenVLA: An Open-Source Vision-Language-Action Model,'' arXiv:2406.09246, 2024.
\bibitem{pi052025ch5} Physical Intelligence et al., ``$\pi_{0.5}$: A Vision-Language-Action Model with Open-World Generalization,'' arXiv:2504.16054, 2025.
\bibitem{smolvla2025ch5} Hugging Face, ``SmolVLA: Efficient Vision-Language-Action Model trained on Lerobot Community Data,'' Jun. 3, 2025. [Online]. Available: \href{https://huggingface.co/blog/smolvla}{official blog}
\bibitem{genie22024ch5} Google DeepMind, ``Genie 2: A large-scale foundation world model,'' Dec. 4, 2024. [Online]. Available: \href{https://deepmind.google/discover/blog/genie-2-a-large-scale-foundation-world-model/}{official blog}
\bibitem{cosmos2025ch5} NVIDIA, ``NVIDIA Launches Cosmos World Foundation Model Platform to Accelerate Physical AI Development,'' Jan. 6, 2025. [Online]. Available: \href{https://nvidianews.nvidia.com/news/nvidia-launches-cosmos-world-foundation-model-platform-to-accelerate-physical-ai-development}{official newsroom}
\bibitem{helix022026ch5} Figure, ``Introducing Helix 02: Full-Body Autonomy,'' Jan. 27, 2026. [Online]. Available: \href{https://www.figure.ai/news/helix-02}{official page}
\end{thebibliography}
